\documentclass[8pt,aspectratio=169,hyperref={unicode=true}]{beamer}

\usefonttheme{serif}
\usepackage{fontspec}
	\setmainfont{TeX Gyre Heros}
\usepackage{unicode-math}
\usepackage{lualatex-math}
	\setmathfont{TeX Gyre Termes Math}
\usepackage{polyglossia}
\setdefaultlanguage[frenchpart=false]{french}
\setotherlanguage{english}
%\usepackage{microtype}
\usepackage[locale = FR,
            separate-uncertainty,
            multi-part-units = single,
            range-units = single]{siunitx}
	\DeclareSIUnit\an{an}
  \DeclareSIUnit{\octet}{o}
\usepackage{amsmath}
\usepackage{amsfonts}
\usepackage{amssymb}
\usepackage{array}
\usepackage{graphicx}
\graphicspath{{./Figures/}}
\usepackage{booktabs}
\usepackage{tabularx}
\usepackage{multirow}
\usepackage{multicol}
    \newcolumntype{L}{>{\raggedright\arraybackslash}X}
    \newcolumntype{R}{>{\raggedleft\arraybackslash}X}
\usepackage{tikz}
\usetikzlibrary{graphs, graphdrawing} \usegdlibrary{trees}
\usetikzlibrary{overlay-beamer-styles}
\usepackage{subcaption}
\usepackage[]{animate}
\usepackage{float}
\usepackage{csquotes}

\usetheme[secheader
         ]{Boadilla}
\usecolortheme{seagull}
\setbeamertemplate{enumerate items}[default]
\setbeamertemplate{itemize items}{-}
\setbeamertemplate{navigation symbols}{}
\setbeamertemplate{bibliography item}{}
\setbeamerfont{framesubtitle}{size=\large}
\setbeamertemplate{section in toc}[sections numbered]
%\setbeamertemplate{subsection in toc}[subsections numbered]

\title[Anticipez les besoins en consommation électrique de bâtiments]
{Projet 4 : Anticipez les besoins en consommation électrique de bâtiments}
\author[Lancelot \textsc{Leclercq}]{Lancelot \textsc{Leclercq}} 
\institute[]{}
\date[]{\small{15 décembre 2021}}

\AtBeginSection[]{
  \begin{frame}
  \vfill
  \centering
    \usebeamerfont{title}\insertsectionhead\par%
  \vfill
  \end{frame}
}

\begin{document}
\setbeamercolor{background canvas}{bg=gray!20}
\begin{frame}[plain]
  \titlepage
\end{frame}

\begin{frame}{Sommaire}
  \Large
  \begin{columns}
    \begin{column}{.7\textwidth}
      \tableofcontents[hideallsubsections]
    \end{column}
  \end{columns}
\end{frame}


\section{Introduction}
\subsection{Problématique}
\begin{frame}{\insertsubsection}
  \begin{columns}
    \begin{column}{.6\textwidth}
      \begin{itemize}
        \item Objectif de la ville de Seattle : atteindre la neutralité en émissions
              de carbone
        \item[]
        \item La ville s'intéresse aux émissions des batiments non destinés
              à l'habitation
        \item[]
        \item Pour cela des relevés de consommation ont été réalisés mais ils sont
              couteux à obtenir
        \item[]
        \item Est-il possible de prédire les émissions et de la consommation d'énergie
              pour des batiments pour lesquels les relevés n'ont pas été réalisé à partir
              des relevés déjà obtenus
      \end{itemize}
    \end{column}
    \begin{column}{.4\textwidth}
      \begin{figure}
        \includegraphics[width=.8\textwidth]{./Seattle_logo_landscape_blue-black.pdf}
      \end{figure}
    \end{column}
  \end{columns}
\end{frame}

\subsection{Jeu de données}
\begin{frame}{\insertsubsection}
  \begin{itemize}
    \item Base de données issue de l'initiative de la ville de Seattle de proposer ses
          données en accès libre (Open Data)
    \item[]
    \item Données concernant les batiments de la ville, caractérise :
          \begin{itemize}
            \item le type,
            \item la surface,
            \item le nombre d'étages,
            \item la consomation énergétique,
            \item les émissions de carbone,
            \item $\vdots$
          \end{itemize}
    \item[]
    \item Données des années 2015 et 2016
  \end{itemize}
\end{frame}

\section{Nettoyage du jeu de données}
\subsection{Correction et selection des données}
\begin{frame}{\insertsection : \insertsubsection}
  \begin{columns}[t]
    \begin{column}{.5\textwidth}
      \includegraphics[width=\textwidth]{DataNbFull.pdf}
      \begin{itemize}
        \item Nettoyage des valeurs négative pour la surface des batiments/parkings,
              la consommation et les émissions
        \item Correction du nombre de d'étages aberrant pour certains batiments
        \item Correction du nombre de batiment nul par 1
      \end{itemize}
    \end{column}
    \begin{column}{.5\textwidth}
      \includegraphics[width=\textwidth]{DataNbDrop.pdf}
      \begin{itemize}
        \item Conservation des variables ayant plus de 50\% de données
        \item Suppressions des variables étant des relevés afin de voir si notre modèle
              peut s'en passer
      \end{itemize}
    \end{column}
  \end{columns}
\end{frame}

\subsection{Selections des variables}
\subsubsection{RFE et matrice de corrélation}
\begin{frame}{\insertsection : \insertsubsection}{\insertsubsubsection}
  \begin{columns}[t]
    \begin{column}{.5\textwidth}
      \centering
      Variables pertinentes pour les émissions
      \includegraphics[width=\textwidth]{RFECVEmissions.pdf}
    \end{column}
    \begin{column}{.5\textwidth}
      \centering
      Variables pertinentes pour la consommation
      \includegraphics[width=\textwidth]{RFECVConso.pdf}
    \end{column}
  \end{columns}
  \begin{itemize}
    \item Selection des variables les plus pertinentes par elimination recursive des variables (RFE)
    \item Réduction efficace pour les émissions
    \item Pas de réel changement de RMSE pour la consommation
  \end{itemize}
\end{frame}

\begin{frame}{\insertsection : \insertsubsection}{\insertsubsubsection}
  \begin{columns}[t]
    \begin{column}{.33\textwidth}
      \centering
      Variables pertinentes pour les émissions
      \includegraphics[width=\textwidth]{HeatmapUsedNumEmissions.pdf}
    \end{column}
    \begin{column}{.33\textwidth}
      \begin{itemize}
        \item Observation des résultats de RFE par les matrices de corrélation
        \item[]
        \item Les variables les plus corrélées sont communes aux deux sélection
        \item[]
        \item Conservation de 6 variables jugées pertinentes
      \end{itemize}
    \end{column}
    \begin{column}{.33\textwidth}
      \centering
      Variables pertinentes pour la consommation
      \includegraphics[width=\textwidth]{HeatmapUsedNumConso.pdf}
    \end{column}
  \end{columns}
\end{frame}

\subsubsection{PCA}
\begin{frame}[t]{\insertsection : \insertsubsection}{\insertsubsubsection}
  \begin{columns}
    \begin{column}{.33\textwidth}
      \includegraphics[width=\textwidth]{ScreePlot.pdf}
      \begin{itemize}
        \item Le graphique de la variance expliquée cumulée nous montre que
              99\% de la matrice est exliquée avec 5 variables
        \item Les quatres variables les plus corrélées se retrouvent sur l'axe F1
        \item L'EnergyStar score semble avoir une certaine importance car il explique
              une grande partie de l'axe F3
      \end{itemize}
    \end{column}
    \begin{column}{.33\textwidth}
      \centering
      \includegraphics[width=\textwidth]{PCAF1F2.pdf}
    \end{column}
    \begin{column}{.33\textwidth}
      \centering
      \includegraphics[width=\textwidth]{PCAF1F3.pdf}
    \end{column}
  \end{columns}
\end{frame}

\section{Étapes des modélisations}
\begin{frame}{\insertsection}
  Afin de comparer les différents modèles
  \begin{itemize}
    \item split commun à chaque modèle (varie selon la variable modélisée)
    \item boucle pour chaque modèle
          \begin{itemize}
            \item création d'un pipeline : scaling et fit du modèle
                  \begin{itemize}
                    \item scaling par RobustScaler car plus résistant aux valeurs aberrantes
                          selon la documentation
                  \end{itemize}
          \end{itemize}
    \item la boucle retourne :
          \begin{itemize}
            \item la RMSE en fonction du paramètre le plus évolutif
            \item le(s) meilleur(s) paramètre(s)
            \item le R², la RMSE, la MAE (mean absolute error) et le temps de calcul du modèle
          \end{itemize}
  \end{itemize}
  \vspace{.5cm}
  \tikz [nodes={draw}]
  \graph [tree layout, level distance=4cm] {
  train test split [grow=right] ->
  {RobustScaler -> "Modèle.fit"} -> {"R², RMSE, MAE et temps de calcul",
  meilleur paramètre, "figure RMSE", 
  figure variable test vs variable prédite};
  boucle [draw] // {RobustScaler, "Modèle.fit", "R², RMSE, MAE et temps de calcul",
  meilleur paramètre, "figure RMSE",
  figure variable test vs variable prédite};
  pipeline [draw] // {RobustScaler, "Modèle.fit"};
  };

\end{frame}

\section[Modélisation émissions]{Modélisation des émissions de carbone}
\subsection{Modèle Ridge}
\begin{frame}{\insertsubsection}
  \begin{columns}[t]
    \begin{column}{.33\textwidth}
      \centering Variable non modifiée
      \includegraphics[width=.97\textwidth]{EmissionsGraphRMSERidge.pdf}
      \includegraphics[width=.97\textwidth]{EmissionsTestvsPredRidge.pdf}
    \end{column}
    \begin{column}{.33\textwidth}
      $\Longleftarrow$
      \scriptsize
      {\centering
        \input{./Tableaux/EmissionsScoreRidge.tex}}
      \input{./Tableaux/EmissionsBestParamsRidge.tex}

      \normalsize
      $\Longleftarrow$

      \raggedleft{$\Longrightarrow$}
      \scriptsize
      {\centering
        \input{./Tableaux/EmissionsScoreRidge_log.tex}}
      \input{./Tableaux/EmissionsBestParamsRidge_log.tex}

      \normalsize
      \raggedleft{$\Longrightarrow$}
    \end{column}
    \begin{column}{.33\textwidth}
      \centering Variable au log
      \includegraphics[width=.96\textwidth]{EmissionsGraphRMSERidge_log.pdf}
      \includegraphics[width=.96\textwidth]{EmissionsTestvsPredRidge_log.pdf}
    \end{column}
  \end{columns}
\end{frame}

\subsection{Modèle Lasso}
\begin{frame}{\insertsubsection}
  \begin{columns}[t]
    \begin{column}{.33\textwidth}
      \centering Variable non modifiée
      \includegraphics[width=.97\textwidth]{EmissionsGraphRMSELasso.pdf}
      \includegraphics[width=.97\textwidth]{EmissionsTestvsPredLasso.pdf}
    \end{column}
    \begin{column}{.33\textwidth}
      $\Longleftarrow$
      \scriptsize
      {\centering
        \input{./Tableaux/EmissionsScoreLasso.tex}}
      \input{./Tableaux/EmissionsBestParamsLasso.tex}

      \normalsize
      $\Longleftarrow$

      \raggedleft{$\Longrightarrow$}
      \scriptsize
      {\centering
        \input{./Tableaux/EmissionsScoreLasso_log.tex}}
      \input{./Tableaux/EmissionsBestParamsLasso_log.tex}

      \normalsize
      \raggedleft{$\Longrightarrow$}
    \end{column}
    \begin{column}{.33\textwidth}
      \centering Variable au log
      \includegraphics[width=.96\textwidth]{EmissionsGraphRMSELasso_log.pdf}
      \includegraphics[width=.96\textwidth]{EmissionsTestvsPredLasso_log.pdf}
    \end{column}
  \end{columns}
\end{frame}

\subsection{Modèle ElasticNet}
\begin{frame}{\insertsubsection}
  \begin{columns}[t]
    \begin{column}{.33\textwidth}
      \centering Variable non modifiée
      \includegraphics[width=.97\textwidth]{EmissionsGraphRMSEElasticNet.pdf}
      \includegraphics[width=.97\textwidth]{EmissionsTestvsPredElasticNet.pdf}
    \end{column}
    \begin{column}{.33\textwidth}
      $\Longleftarrow$
      \scriptsize
      {\centering
        \input{./Tableaux/EmissionsScoreElasticNet.tex}}
      \input{./Tableaux/EmissionsBestParamsElasticNet.tex}

      \normalsize
      $\Longleftarrow$

      \raggedleft{$\Longrightarrow$}
      \scriptsize
      {\centering
        \input{./Tableaux/EmissionsScoreElasticNet_log.tex}}
      \input{./Tableaux/EmissionsBestParamsElasticNet_log.tex}

      \normalsize
      \raggedleft{$\Longrightarrow$}
    \end{column}
    \begin{column}{.33\textwidth}
      \centering Variable au log
      \includegraphics[width=.96\textwidth]{EmissionsGraphRMSEElasticNet_log.pdf}
      \includegraphics[width=.96\textwidth]{EmissionsTestvsPredElasticNet_log.pdf}
    \end{column}
  \end{columns}
\end{frame}

\subsection{Modèle kNeighborsRegressor}
\begin{frame}{\insertsubsection}
  \begin{columns}[t]
    \begin{column}{.33\textwidth}
      \centering Variable non modifiée
      \includegraphics[width=.97\textwidth]{EmissionsGraphRMSEKNeighborsRegressor.pdf}
      \includegraphics[width=.97\textwidth]{EmissionsTestvsPredKNeighborsRegressor.pdf}
    \end{column}
    \begin{column}{.33\textwidth}
      $\Longleftarrow$
      \scriptsize
      {\centering
        \input{./Tableaux/EmissionsScoreKNeighborsRegressor.tex}}
      \input{./Tableaux/EmissionsBestParamsKNeighborsRegressor.tex}

      \normalsize
      $\Longleftarrow$

      \raggedleft{$\Longrightarrow$}
      \scriptsize
      {\centering
        \input{./Tableaux/EmissionsScoreKNeighborsRegressor_log.tex}}
      \input{./Tableaux/EmissionsBestParamsKNeighborsRegressor_log.tex}

      \normalsize
      \raggedleft{$\Longrightarrow$}
    \end{column}
    \begin{column}{.33\textwidth}
      \centering Variable au log
      \includegraphics[width=.96\textwidth]{EmissionsGraphRMSEKNeighborsRegressor_log.pdf}
      \includegraphics[width=.96\textwidth]{EmissionsTestvsPredKNeighborsRegressor_log.pdf}
    \end{column}
  \end{columns}
\end{frame}

\subsection{Modèle RandomForestRegressor}
\begin{frame}{\insertsubsection}
  \begin{columns}[t]
    \begin{column}{.33\textwidth}
      \centering Variable non modifiée
      \includegraphics[width=.97\textwidth]{EmissionsGraphRMSERandomForestRegressor.pdf}
      \includegraphics[width=.97\textwidth]{EmissionsTestvsPredRandomForestRegressor.pdf}
    \end{column}
    \begin{column}{.33\textwidth}
      $\Longleftarrow$
      \scriptsize
      {\centering
        \input{./Tableaux/EmissionsScoreRandomForestRegressor.tex}}
      \input{./Tableaux/EmissionsBestParamsRandomForestRegressor.tex}

      \normalsize
      $\Longleftarrow$

      \raggedleft{$\Longrightarrow$}
      \scriptsize
      {\centering
        \input{./Tableaux/EmissionsScoreRandomForestRegressor_log.tex}}
      \input{./Tableaux/EmissionsBestParamsRandomForestRegressor_log.tex}

      \normalsize
      \raggedleft{$\Longrightarrow$}
    \end{column}
    \begin{column}{.33\textwidth}
      \centering Variable au log
      \includegraphics[width=.96\textwidth]{EmissionsGraphRMSERandomForestRegressor_log.pdf}
      \includegraphics[width=.96\textwidth]{EmissionsTestvsPredRandomForestRegressor_log.pdf}
    \end{column}
  \end{columns}
\end{frame}

\subsection{Modèle AdaBoostRegressor}
\begin{frame}{\insertsubsection}
  \begin{columns}[t]
    \begin{column}{.33\textwidth}
      \centering Variable non modifiée
      \includegraphics[width=.97\textwidth]{EmissionsGraphRMSEAdaBoostRegressor.pdf}
      \includegraphics[width=.97\textwidth]{EmissionsTestvsPredAdaBoostRegressor.pdf}
    \end{column}
    \begin{column}{.33\textwidth}
      $\Longleftarrow$
      \scriptsize
      {\centering
        \input{./Tableaux/EmissionsScoreAdaBoostRegressor.tex}}
      \input{./Tableaux/EmissionsBestParamsAdaBoostRegressor.tex}

      \normalsize
      $\Longleftarrow$

      \raggedleft{$\Longrightarrow$}
      \scriptsize
      {\centering
        \input{./Tableaux/EmissionsScoreAdaBoostRegressor_log.tex}}
      \input{./Tableaux/EmissionsBestParamsAdaBoostRegressor_log.tex}

      \normalsize
      \raggedleft{$\Longrightarrow$}
    \end{column}
    \begin{column}{.33\textwidth}
      \centering Variable au log
      \includegraphics[width=.96\textwidth]{EmissionsGraphRMSEAdaBoostRegressor_log.pdf}
      \includegraphics[width=.96\textwidth]{EmissionsTestvsPredAdaBoostRegressor_log.pdf}
    \end{column}
  \end{columns}
\end{frame}

\subsection{Modèle GradientBoostingRegressor}
\begin{frame}{\insertsubsection}
  \begin{columns}[t]
    \begin{column}{.33\textwidth}
      \centering Variable non modifiée
      \includegraphics[width=.97\textwidth]{EmissionsGraphRMSEGradientBoostingRegressor.pdf}
      \includegraphics[width=.97\textwidth]{EmissionsTestvsPredGradientBoostingRegressor.pdf}
    \end{column}
    \begin{column}{.33\textwidth}
      $\Longleftarrow$
      \scriptsize
      {\centering
        \input{./Tableaux/EmissionsScoreGradientBoostingRegressor.tex}}
      \input{./Tableaux/EmissionsBestParamsGradientBoostingRegressor.tex}

      \normalsize
      $\Longleftarrow$

      \raggedleft{$\Longrightarrow$}
      \scriptsize
      {\centering
        \input{./Tableaux/EmissionsScoreGradientBoostingRegressor_log.tex}}
      \input{./Tableaux/EmissionsBestParamsGradientBoostingRegressor_log.tex}

      \normalsize
      \raggedleft{$\Longrightarrow$}
    \end{column}
    \begin{column}{.33\textwidth}
      \centering Variable au log
      \includegraphics[width=.96\textwidth]{EmissionsGraphRMSEGradientBoostingRegressor_log.pdf}
      \includegraphics[width=.96\textwidth]{EmissionsTestvsPredGradientBoostingRegressor_log.pdf}
    \end{column}
  \end{columns}
\end{frame}

\subsection{Comparaison des résultats selon que la variable est au log ou non}
\begin{frame}{\insertsubsection}
  \begin{columns}
    \begin{column}{.5\textwidth}
      \centering
      \includegraphics[width=.95\textwidth]{EmissionsCompareScores.pdf}
    \end{column}
    \begin{column}{.5\textwidth}
      \begin{itemize}
        \item
      \end{itemize}
    \end{column}
  \end{columns}
\end{frame}

\subsection{Influence de l'EnergyStar score sur la prédiction des Émissions}
\begin{frame}{\insertsubsection}
  \begin{columns}
    \begin{column}{.6\textwidth}
      \centering
      \includegraphics[width=\textwidth]{EmissionsCompareScoresESS.pdf}
    \end{column}
    \begin{column}{.4\textwidth}
      \begin{itemize}
        \item
      \end{itemize}
    \end{column}
  \end{columns}
\end{frame}

\section[Modélisation consommation]{Modélisation de la consommation énergétique}
\subsection{Modèle Ridge}
\begin{frame}{\insertsubsection}
  \begin{columns}[t]
    \begin{column}{.33\textwidth}
      \centering Variable non modifiée
      \includegraphics[width=.97\textwidth]{ConsoGraphRMSERidge.pdf}
      \includegraphics[width=.97\textwidth]{ConsoTestvsPredRidge.pdf}
    \end{column}
    \begin{column}{.33\textwidth}
      $\Longleftarrow$
      \scriptsize
      {\centering
        \input{./Tableaux/ConsoScoreRidge.tex}}
      \input{./Tableaux/ConsoBestParamsRidge.tex}

      \normalsize
      $\Longleftarrow$

      \raggedleft{$\Longrightarrow$}
      \scriptsize
      {\centering
        \input{./Tableaux/ConsoScoreRidge_log.tex}}
      \input{./Tableaux/ConsoBestParamsRidge_log.tex}

      \normalsize
      \raggedleft{$\Longrightarrow$}
    \end{column}
    \begin{column}{.33\textwidth}
      \centering Variable au log
      \includegraphics[width=.96\textwidth]{ConsoGraphRMSERidge_log.pdf}
      \includegraphics[width=.96\textwidth]{ConsoTestvsPredRidge_log.pdf}
    \end{column}
  \end{columns}
\end{frame}

\subsection{Modèle Lasso}
\begin{frame}{\insertsubsection}
  \begin{columns}[t]
    \begin{column}{.33\textwidth}
      \centering Variable non modifiée
      \includegraphics[width=.97\textwidth]{ConsoGraphRMSELasso.pdf}
      \includegraphics[width=.97\textwidth]{ConsoTestvsPredLasso.pdf}
    \end{column}
    \begin{column}{.33\textwidth}
      $\Longleftarrow$
      \scriptsize
      {\centering
        \input{./Tableaux/ConsoScoreLasso.tex}}
      \input{./Tableaux/ConsoBestParamsLasso.tex}

      \normalsize
      $\Longleftarrow$

      \raggedleft{$\Longrightarrow$}
      \scriptsize
      {\centering
        \input{./Tableaux/ConsoScoreLasso_log.tex}}
      \input{./Tableaux/ConsoBestParamsLasso_log.tex}

      \normalsize
      \raggedleft{$\Longrightarrow$}
    \end{column}
    \begin{column}{.33\textwidth}
      \centering Variable au log
      \includegraphics[width=.96\textwidth]{ConsoGraphRMSELasso_log.pdf}
      \includegraphics[width=.96\textwidth]{ConsoTestvsPredLasso_log.pdf}
    \end{column}
  \end{columns}
\end{frame}

\subsection{Modèle ElasticNet}
\begin{frame}{\insertsubsection}
  \begin{columns}[t]
    \begin{column}{.33\textwidth}
      \centering Variable non modifiée
      \includegraphics[width=.97\textwidth]{ConsoGraphRMSEElasticNet.pdf}
      \includegraphics[width=.97\textwidth]{ConsoTestvsPredElasticNet.pdf}
    \end{column}
    \begin{column}{.33\textwidth}
      $\Longleftarrow$
      \scriptsize
      {\centering
        \input{./Tableaux/ConsoScoreElasticNet.tex}}
      \input{./Tableaux/ConsoBestParamsElasticNet.tex}

      \normalsize
      $\Longleftarrow$

      \raggedleft{$\Longrightarrow$}
      \scriptsize
      {\centering
        \input{./Tableaux/ConsoScoreElasticNet_log.tex}}
      \input{./Tableaux/ConsoBestParamsElasticNet_log.tex}

      \normalsize
      \raggedleft{$\Longrightarrow$}
    \end{column}
    \begin{column}{.33\textwidth}
      \centering Variable au log
      \includegraphics[width=.96\textwidth]{ConsoGraphRMSEElasticNet_log.pdf}
      \includegraphics[width=.96\textwidth]{ConsoTestvsPredElasticNet_log.pdf}
    \end{column}
  \end{columns}
\end{frame}

\subsection{Modèle kNeighborsRegressor}
\begin{frame}{\insertsubsection}
  \begin{columns}[t]
    \begin{column}{.33\textwidth}
      \centering Variable non modifiée
      \includegraphics[width=.97\textwidth]{ConsoGraphRMSEKNeighborsRegressor.pdf}
      \includegraphics[width=.97\textwidth]{ConsoTestvsPredKNeighborsRegressor.pdf}
    \end{column}
    \begin{column}{.33\textwidth}
      $\Longleftarrow$
      \scriptsize
      {\centering
        \input{./Tableaux/ConsoScoreKNeighborsRegressor.tex}}
      \input{./Tableaux/ConsoBestParamsKNeighborsRegressor.tex}

      \normalsize
      $\Longleftarrow$

      \raggedleft{$\Longrightarrow$}
      \scriptsize
      {\centering
        \input{./Tableaux/ConsoScoreKNeighborsRegressor_log.tex}}
      \input{./Tableaux/ConsoBestParamsKNeighborsRegressor_log.tex}

      \normalsize
      \raggedleft{$\Longrightarrow$}
    \end{column}
    \begin{column}{.33\textwidth}
      \centering Variable au log
      \includegraphics[width=.96\textwidth]{ConsoGraphRMSEKNeighborsRegressor_log.pdf}
      \includegraphics[width=.96\textwidth]{ConsoTestvsPredKNeighborsRegressor_log.pdf}
    \end{column}
  \end{columns}
\end{frame}

\subsection{Modèle RandomForestRegressor}
\begin{frame}{\insertsubsection}
  \begin{columns}[t]
    \begin{column}{.33\textwidth}
      \centering Variable non modifiée
      \includegraphics[width=.97\textwidth]{ConsoGraphRMSERandomForestRegressor.pdf}
      \includegraphics[width=.97\textwidth]{ConsoTestvsPredRandomForestRegressor.pdf}
    \end{column}
    \begin{column}{.33\textwidth}
      $\Longleftarrow$
      \scriptsize
      {\centering
        \input{./Tableaux/ConsoScoreRandomForestRegressor.tex}}
      \input{./Tableaux/ConsoBestParamsRandomForestRegressor.tex}

      \normalsize
      $\Longleftarrow$

      \raggedleft{$\Longrightarrow$}
      \scriptsize
      {\centering
        \input{./Tableaux/ConsoScoreRandomForestRegressor_log.tex}}
      \input{./Tableaux/ConsoBestParamsRandomForestRegressor_log.tex}

      \normalsize
      \raggedleft{$\Longrightarrow$}
    \end{column}
    \begin{column}{.33\textwidth}
      \centering Variable au log
      \includegraphics[width=.96\textwidth]{ConsoGraphRMSERandomForestRegressor_log.pdf}
      \includegraphics[width=.96\textwidth]{ConsoTestvsPredRandomForestRegressor_log.pdf}
    \end{column}
  \end{columns}
\end{frame}

\subsection{Modèle AdaBoostRegressor}
\begin{frame}{\insertsubsection}
  \begin{columns}[t]
    \begin{column}{.33\textwidth}
      \centering Variable non modifiée
      \includegraphics[width=.97\textwidth]{ConsoGraphRMSEAdaBoostRegressor.pdf}
      \includegraphics[width=.97\textwidth]{ConsoTestvsPredAdaBoostRegressor.pdf}
    \end{column}
    \begin{column}{.33\textwidth}
      $\Longleftarrow$
      \scriptsize
      {\centering
        \input{./Tableaux/ConsoScoreAdaBoostRegressor.tex}}
      \input{./Tableaux/ConsoBestParamsAdaBoostRegressor.tex}

      \normalsize
      $\Longleftarrow$

      \raggedleft{$\Longrightarrow$}
      \scriptsize
      {\centering
        \input{./Tableaux/ConsoScoreAdaBoostRegressor_log.tex}}
      \input{./Tableaux/ConsoBestParamsAdaBoostRegressor_log.tex}

      \normalsize
      \raggedleft{$\Longrightarrow$}
    \end{column}
    \begin{column}{.33\textwidth}
      \centering Variable au log
      \includegraphics[width=.96\textwidth]{ConsoGraphRMSEAdaBoostRegressor_log.pdf}
      \includegraphics[width=.96\textwidth]{ConsoTestvsPredAdaBoostRegressor_log.pdf}
    \end{column}
  \end{columns}
\end{frame}

\subsection{Modèle GradientBoostingRegressor}
\begin{frame}{\insertsubsection}
  \begin{columns}[t]
    \begin{column}{.33\textwidth}
      \centering Variable non modifiée
      \includegraphics[width=.97\textwidth]{ConsoGraphRMSEGradientBoostingRegressor.pdf}
      \includegraphics[width=.97\textwidth]{ConsoTestvsPredGradientBoostingRegressor.pdf}
    \end{column}
    \begin{column}{.33\textwidth}
      $\Longleftarrow$
      \scriptsize
      {\centering
        \input{./Tableaux/ConsoScoreGradientBoostingRegressor.tex}}
      \input{./Tableaux/ConsoBestParamsGradientBoostingRegressor.tex}

      \normalsize
      $\Longleftarrow$

      \raggedleft{$\Longrightarrow$}
      \scriptsize
      {\centering
        \input{./Tableaux/ConsoScoreGradientBoostingRegressor_log.tex}}
      \input{./Tableaux/ConsoBestParamsGradientBoostingRegressor_log.tex}

      \normalsize
      \raggedleft{$\Longrightarrow$}
    \end{column}
    \begin{column}{.33\textwidth}
      \centering Variable au log
      \includegraphics[width=.96\textwidth]{ConsoGraphRMSEGradientBoostingRegressor_log.pdf}
      \includegraphics[width=.96\textwidth]{ConsoTestvsPredGradientBoostingRegressor_log.pdf}
    \end{column}
  \end{columns}
\end{frame}

\subsection{Comparaison des résultats selon que la variable est au log ou non}
\begin{frame}{\insertsubsection}
  \begin{columns}
    \begin{column}{.5\textwidth}
      \centering
      \includegraphics[width=.95\textwidth]{ConsoCompareScores.pdf}
    \end{column}
    \begin{column}{.5\textwidth}
      \begin{itemize}
        \item
      \end{itemize}
    \end{column}
  \end{columns}
\end{frame}

\section{Conclusion}
\begin{frame}{\insertsection}

\end{frame}

\end{document}